\documentclass[12pt]{article}
\usepackage[margin=0.5in,top=0.7in,bottom=0.45in]{geometry}
\usepackage{lmodern}
\usepackage{microtype}
\usepackage{fancyhdr}
\usepackage{titlesec}
\usepackage{enumitem}
\usepackage{lastpage}
\usepackage[hidelinks]{hyperref}

\setlength{\parindent}{0pt}
\setlength{\parskip}{0.1em}
\setlength{\headheight}{26pt}
\titleformat{\section}{\normalsize\bfseries\scshape}{}{0pt}{}
\titlespacing*{\section}{0pt}{0.28em}{0.1em}
\pagestyle{fancy}
\fancyhf{}
\fancyhead[C]{\footnotesize Student Name: \hspace{2.3in} \quad Assignment ID: U1L6LE \quad Page \thepage\ of \pageref{LastPage}}
\renewcommand{\headrulewidth}{0.5pt}

\newcommand{\answerbox}[2]{%
\vspace{0.04em}
\noindent\textbf{#1}\par
\vspace{0.01em}
\noindent\fbox{%
\begin{minipage}[t][#2]{\dimexpr\textwidth-2\fboxsep-2\fboxrule}
\rule{\textwidth}{0pt}
\vspace{#2}
\end{minipage}}
\par\vspace{0.08em}}

\begin{document}

\begin{center}
{\LARGE\bfseries Week 6 Lit Example Worksheet}\par\vspace{0.02em}
{\large\scshape Julius Caesar: Valid Form, Unproved Premise}\par\vspace{0.05em}
\rule{0.78\textwidth}{0.5pt}
\end{center}

\textbf{Purpose:} Use the Lit Example Reader. Apply Week 6's form test and fact test to Brutus and Antony separately; do not collapse validity, truth, and persuasion.

\section*{Part 1: Form Test}

\textbf{Part 1 Q1:} Reconstruct Brutus's argument in this form: All M are P; All S are M; therefore All S are P. Name M, S, and P.

\answerbox{Part 1 Q1 ANSWER BOX}{1.2in}

\textbf{Part 1 Q2:} Apply the form test only. If Brutus's two premises were true, would the conclusion follow? Explain.

\answerbox{Part 1 Q2 ANSWER BOX}{1.0in}

\section*{Part 2: Fact Test}

\textbf{Part 2 Q1:} Which premise in Brutus's argument most needs factual proof? Explain why it is the pressure-point premise.

\answerbox{Part 2 Q1 ANSWER BOX}{1.1in}

\textbf{Part 2 Q2:} What evidence would Brutus need in order to make that premise strong enough for a careful listener?

\answerbox{Part 2 Q2 ANSWER BOX}{1.1in}

\newpage

\section*{Part 3: Antony's Evidence}

\textbf{Part 3 Q1:} List Antony's three main pieces of evidence against Brutus's ambition premise.

\answerbox{Part 3 Q1 ANSWER BOX}{1.15in}

\textbf{Part 3 Q2:} Does Antony prove that Caesar was not ambitious, or does he mainly make Brutus's premise doubtful? Explain the difference.

\answerbox{Part 3 Q2 ANSWER BOX}{1.4in}

\section*{Part 4: Final Diagnosis}

\textbf{Part 4 Q1:} Diagnose Brutus's argument as valid, invalid, sound, unsound, or not yet proved. You may use more than one label if you explain the difference.

\answerbox{Part 4 Q1 ANSWER BOX}{1.4in}

\textbf{Part 4 Q2:} In one paragraph, explain why ``valid'' does not mean ``true'' using Brutus and Antony as your example.

\answerbox{Part 4 Q2 ANSWER BOX}{1.35in}

\end{document}
